\section{Conclusion}
\label{sec:conclusion}

We trained a 4.88M-parameter ResNet from scratch on CIFAR-10 and reached 97.84\%
clean validation accuracy (97.92\% with 10-view TTA). The team's selected
submissions scored 0.911 and 0.904 public, with a private score of 0.880 and
13th place. The
recipe is ordinary in its ingredients---long cosine training, SAM, EMA,
TrivialAugment, Cutout, mixup/CutMix, and three-teacher KD---but the combination
is strong at this parameter budget, and the student beats its largest teacher.

The project also produced a useful catalog of negative results. RandAugment,
AutoAugment-CIFAR, geometry-only and affine+color augmentation, ASAM/SWA
fine-tuning, logit-standardized KD, a two-model parameter split, multi-crop TTA,
test-time BN adaptation, Sinkhorn balancing, and seed ensembling all failed to
improve the metric in our hands. One-view AugMix was the exception: clean-neutral and
consistently better on corrupted and photometric proxies, but not shipped because
its KD confirmation finished after submission selection.

The unresolved question is the roughly ten-point gap between clean validation and
the private leaderboard. We do not have the final evaluation images or labels,
so we can describe the gap but not explain it. We hope the
negative results and the explicit separation between proxy evidence and
leaderboard evidence are useful to others working under similar constraints.

\paragraph{Reproducibility.}
Every run writes \texttt{args.json}, a per-epoch \texttt{log.csv}, a
\texttt{results.json}, and a resumable \texttt{last.pt}. The report folder
contains the archived logs, the fixed figure-generation script, and the AAAI
style file; every figure can be rebuilt without the competition data or any
checkpoint. The submitted models are 4.88M and 4.92M parameters, trained from
scratch with no external data.
